\documentclass[11pt]{amsart}

\usepackage[T1]{fontenc}
\usepackage{lmodern}
\usepackage{microtype}
\usepackage{amsmath,amssymb,amsthm}
\usepackage[hidelinks]{hyperref}

\hypersetup{
  pdftitle={An Infinite Family of Counterexamples to an Atomic-Weight Conjecture of Dargad and Larsson}
}

\newtheorem{theorem}{Theorem}
\newtheorem{corollary}[theorem]{Corollary}
\theoremstyle{remark}
\newtheorem*{remark}{Remark}

\newcommand{\TS}{\mathcal{T}}
\newcommand{\Up}{\mathord{\uparrow}}
\newcommand{\Down}{\mathord{\downarrow}}
\newcommand{\aw}{\operatorname{aw}}

\title[Counterexamples to an atomic-weight conjecture]
{An Infinite Family of Counterexamples to an Atomic-Weight Conjecture
of Dargad and Larsson}

\date{}

\subjclass[2020]{91A46, 91A05}
\keywords{partizan subtraction game, truncated support, atomic weight, counterexample}

\begin{document}

\begin{abstract}
Dargad and Larsson conjectured an atomic-weight formula for
truncated-support partizan subtraction games. We prove that, for every
integer $t\geq1$,
\[
  \TS_t(2t+4;1,2t+1)=\Down.
\]
Its atomic weight is $-1$, whereas Conjecture~7.1 in version~1 of their
preprint predicts $0$. In these instances $a=1$, and the formula fails
already at the smallest heap size it governs. The equality
$\tau=(b-a)/2$, which these instances also satisfy, is not itself the
obstruction: the formula does hold at that boundary in instances
tabulated by Dargad and Larsson. Their theorem for $\tau<(b-a)/2$ is
unaffected.
\end{abstract}

\maketitle

\section{The counterexamples}

In the truncated-support game $\TS_\tau(n;a,b)$, Left may subtract any
integer in $[a]=\{1,\dots,a\}$ from a heap of size $n$, while Right may
subtract any integer in $\{\tau+1,\dots,b\}$
\cite[Definition~5.1]{DL}. Play is normal: a player with no legal move
loses \cite{ANW}. In version~1 of their preprint, Dargad and
Larsson conjectured that, for $a<b$, $\tau<b-a$, and
$n\geq R_{\widehat\beta}$,
\begin{equation}\label{eq:conjecture}
  \aw\bigl(\TS_\tau(n;a,b)\bigr)
  =-
  \left\lfloor
    \frac{n-R_{\widehat\beta}}{a}
  \right\rfloor,
\end{equation}
where
\[
  \widehat\beta=
  \left\lfloor\frac{b-a}{b-a-\tau}\right\rfloor,
  \qquad
  R_{\widehat\beta}=\widehat\beta(a+\tau+1)
\]
\cite[Equations~(5), (13), and Conjecture~7.1]{DL}.

We use the standard notation
\[
  *=\{0\mid0\},
  \qquad
  \Up=\{0\mid *\},
  \qquad
  \Down=\{*\mid0\}=-\Up.
\]
We also use the standard fact that a short game equals $0$ precisely
when the second player wins, whether the second player is Left or Right
\cite{ANW}.

\begin{theorem}\label{thm:family}
Fix an integer $t\geq1$ and write
\[
  H_n=\TS_t(n;1,2t+1).
\]
Then
\[
  H_{2t+2}=0,
  \qquad
  H_{2t+3}=*,
  \qquad
  H_{2t+4}=\Down.
\]
\end{theorem}

\begin{proof}
Here Left removes one token, while Right removes any number from
$t+1$ through $2t+1$.

First, $H_{2t+2}=0$. If Left starts, she leaves $2t+1$ tokens and Right
removes the entire heap. If Right starts, he leaves a heap of size
$m\in\{1,\dots,t+1\}$; Left moves to $m-1\leq t$, after which Right has
no legal move.

Next consider $H_{2t+3}+*$. If Left first moves in the heap, Right moves
$*\to0$, leaving $H_{2t+2}=0$. If Left instead moves $*\to0$, Right
reduces the heap to $t+2$, after which the heap moves
\[
  t+2\xrightarrow{\mathrm L}t+1\xrightarrow{\mathrm R}0
\]
finish the game.

Suppose instead that Right starts. If he moves $*\to0$, Left reduces the
heap to $2t+2$. If he moves in the heap, he leaves
$m\in\{2,\dots,t+2\}$, and Left moves to
$k=m-1\in\{1,\dots,t+1\}$. A subsequent move $*\to0$ by Right is
answered by $k\to k-1\leq t$, leaving Right without a legal move. A
subsequent heap move by Right is possible only when $k=t+1$, and Left
then moves $*\to0$. Thus $H_{2t+3}+*=0$, so $H_{2t+3}=*$.

Finally consider $H_{2t+4}+\Up$. If Left first moves in the heap, Right
moves $\Up\to *$, leaving $H_{2t+3}+*=0$. If Left instead moves
$\Up\to0$, Right reduces the heap to $t+2$, and the same two heap moves
as above finish the game.

Suppose that Right starts. If he moves $\Up\to *$, Left reduces the heap
to $2t+3$, again leaving $H_{2t+3}+*=0$. If he moves in the heap, he
leaves $m\in\{3,\dots,t+3\}$, and Left moves to
$k=m-1\in\{2,\dots,t+2\}$. If Right next moves in the heap, at most one
token remains; Left moves $\Up\to0$, and Right has no legal move. If
Right instead moves $\Up\to *$, Left reduces the heap to
$k-1\leq t+1$. A move $*\to0$ by Right is then answered by one further
Left heap move, leaving at most $t$ tokens. The only alternative is a
Right heap move, which is possible only from a heap of size $t+1$ and
is answered by $*\to0$. Hence $H_{2t+4}+\Up=0$, and therefore
$H_{2t+4}=-\Up=\Down$.
\end{proof}

\begin{corollary}
Conjecture~7.1 in version~1 of Dargad and Larsson \cite{DL} is false
as stated.
\end{corollary}

\begin{proof}
For any $t\geq1$, take
\[
  (a,b,\tau,n)=(1,2t+1,t,2t+4).
\]
Then $\tau=t<2t=b-a$, and
\[
  \widehat\beta
  =\left\lfloor\frac{2t}{t}\right\rfloor=2,
  \qquad
  R_{\widehat\beta}=2(1+t+1)=2t+4=n.
\]
Formula~\eqref{eq:conjecture} therefore predicts atomic weight $0$.
By Theorem~\ref{thm:family}, however, the game is $\Down$. Since
$\aw(\Up)=1$ and $\aw(-G)=-\aw(G)$,
\[
  \aw(\Down)=-1
\]
\cite[Definition~A.17 and Lemma~A.19(1)]{DL}. This is a contradiction.
\end{proof}

\begin{remark}
Every instance above has $a=1$ and $\widehat\beta=2$, so the failure
occurs at $n=R_{\widehat\beta}$, the smallest heap size to which
\eqref{eq:conjecture} applies. These instances also satisfy
$\tau=(b-a)/2$, but that equality is not itself the obstruction: for
$(a,b)=(3,7)$ and $\tau=2=(b-a)/2$ one has $\widehat\beta=2$ and
$R_{\widehat\beta}=12$, and the atomic weights tabulated in
\cite[Table~4]{DL} agree with \eqref{eq:conjecture} for every tabulated
$n$ with $12\leq n\leq34$. The theorem proved in
\cite[Theorem~1.5]{DL} assumes the strict inequality
$1\leq\tau<(b-a)/2$ and is therefore unaffected. We do not determine
here how far beyond the family of Theorem~\ref{thm:family} the failure
of \eqref{eq:conjecture} extends.
\end{remark}

\begin{thebibliography}{9}

\bibitem{ANW}
M.~H. Albert, R.~J. Nowakowski, and D. Wolfe,
\emph{Lessons in Play: An Introduction to Combinatorial Game Theory},
2nd ed., A K Peters/CRC Press, 2019.

\bibitem{DL}
A. Dargad and U. Larsson,
\emph{Partizan Subtraction with Full and Truncated Support},
\href{https://arxiv.org/abs/2607.27989v1}{arXiv:2607.27989v1}
[math.CO], 2026.

\end{thebibliography}

\end{document}